% Part: many-valued-logic
% Chapter: syntax-and-semantics
% Section: sublogics

\documentclass[../../../include/open-logic-section]{subfiles}

\begin{document}

\olfileid{mvl}{syn}{sub}

\olsection{Logika Mawa Akeh Nilai minangka Sublogika saka~$\LogCL$}

Logika mawa akeh nilai kang lumrah umume ditetepake nganggo matriks
kang fungsi kabenerane, yen diwenehi argumen ing $\{\True, \False\}$,
ngasilake nilai kang padha karo fungsi kabeneran klasik. Mligine,
ing logika kasebut, yen $x \in \{\True, \False\}$, mula
$\tf{\lnot}[\Log L](x) = \tf{\lnot}[\LogCL](x)$; lan yen
$\star$ salah siji saka $\land$, $\lor$, utawa $\lif$,
lan $x, y \in \{\True, \False\}$, mula
$\tf{\star}[\Log L](x,y) = \tf{\star}[\LogCL](x,y)$.
Tegese, fungsi kabeneran kanggo $\lnot$, $\land$, $\lor$,
lan $\lif$ yen diwatesi marang $\{\True,\False\}$
padha persis karo fungsi kabeneran klasik.

\begin{prop}\ollabel{prop:mvl-cl}
  Upamakna basa logika mawa akeh nilai~$\Log L$ mung ngemot
  konektif $\lnot$, $\land$, $\lor$, $\lif$ lan ora ngemot
  tetapan proposisional; $\True, \False \in V$; lan fungsi
  kabenerane nyukupi:
  \begin{enumerate}
    \item\ollabel{prop:not} $\tf{\lnot}[\Log L](x) = \tf{\lnot}[\LogCL](x)$ yen $x =
    \True$ utawa $x = \False$;
    \item\ollabel{prop:land} $\tf{\land}[\Log L](x,y) = \tf{\land}[\LogCL](x,y)$,
    \item\ollabel{prop:lor} $\tf{\lor}[\Log L](x,y) = \tf{\lor}[\LogCL](x,y)$,
    \item\ollabel{prop:lif} $\tf{\lif}[\Log L](x,y) = \tf{\lif}[\LogCL](x,y)$,
    yen $x, y \in \{\True, \False\}$.
  \end{enumerate}
  Mula kanggo saben valuasi $\pAssign v$ menyang~$V$ kang masangake
  nilai Boolean marang saben peubah proposisional, yaiku
  $\pAssign v(p) \in \{\True,\False\}$, kita duwe
  $\pValue v[\Log L](!A) = \pValue v[\LogCL](!A)$.
\end{prop}

\begin{proof}
  Bukti nganggo induksi marang~$!A$.
  \begin{enumerate}
  \item Yen $!A \ident p$ atomik, kita duwe $\pValue v[\Log L](!A) =
  \pAssign v(p) = \pValue
  v[\LogCL](!A)$.
  \item Yen $!A \ident \lnot B$, kita duwe
  \begin{align*}
    \pValue v[\Log L](!A) & = \tf{\lnot}[\Log L](\pValue v[\Log L](!B)) & &\text{miturut \olref[val]{defn:pValue}}\\
    & = \tf{\lnot}[\Log L](\pValue v[\LogCL](!B)) && \text{miturut hipotesis induksi}\\
    & = \tf{\lnot}[\LogCL](\pValue v[\LogCL](!B)) 
&&\text{miturut asumsi \olref{prop:not},}\\
&&&\text{amarga $\pValue v[\LogCL](!B) \in \{\True, \False\}$,}\\
& = \pValue v[\LogCL](!A) &&\text{miturut \olref[val]{defn:pValue}}.
\end{align*}
\item Yen $!A \ident (!B \land !C)$, kita duwe
\begin{align*}
  \pValue v[\Log L](!A) & = \tf{\land}[\Log L](\pValue v[\Log L](!B), \pValue v[\Log L](!C)) & &\text{miturut \olref[val]{defn:pValue}}\\
  & = \tf{\land}[\Log L](\pValue v[\LogCL](!B),\pValue v[\LogCL](!C)) && \text{miturut hipotesis induksi}\\
  & = \tf{\land}[\LogCL](\pValue v[\LogCL](!B),\pValue v[\LogCL](!C)) 
&&\text{miturut asumsi \olref{prop:land},}\\
&&&\text{amarga $\pValue v[\LogCL](!B),\pValue v[\LogCL](!C) \in \{\True, \False\}$,}\\
& = \pValue v[\LogCL](!A) &&\text{miturut \olref[val]{defn:pValue}}.
\end{align*}
\end{enumerate}
Kasus $!A \ident (!B \lor !C)$ lan $!A \ident (!B \lif !C)$
padha carane.
\end{proof}

\begin{cor}
  Yen logika mawa akeh nilai nyukupi syarat ing
  \olref{prop:mvl-cl}, $\True \in V^+$ lan $\False \notin V^+$,
  mula ${\Entails[\Log L]} \subseteq {\Entails[\LogCL]}$;
  tegese, yen $\Gamma
  \Entails[\Log L] !B$, mula $\Gamma \Entails[\LogCL] !B$.
  Mligine, saben tautologi ing $\Log L$ uga tautologi klasik.
\end{cor}

\begin{proof}
  Kita mbuktekake kontraposisine. Upamakna
  $\Gamma \Entails/[\LogCL] !B$. Mula ana
  !!{valuation}~$\pAssign v\colon \PVar \to
  \{\True, \False\}$ kang nyukupi $\pValue v[\LogCL](!A) = \True$
  kanggo saben $!A \in \Gamma$, dene $\pValue v[\LogCL](!B) = \False$.
  Amarga $\True, \False \in V$, !!{valuation}~$\pAssign v$ uga
  !!a{valuation} kanggo~$\Log L$. Miturut \olref{prop:mvl-cl},
  $\pValue v[\Log L](!A) = \True$ kanggo saben $!A \in \Gamma$,
  dene $\pValue v[\Log L](!B) = \False$.
  Amarga $\True \in V^+$ lan $\False \notin V^+$,
  iki ateges $\pSat{v}{\Gamma}[\Log L]$ lan
  $\pSat/{v}{!B}[\Log L]$, yaiku
  $\Gamma \Entails/[\Log L] !B$.
\end{proof}
\end{document}
